% !PRINTED 477
\weberpage{477}
%<<CONT
oder $f_1(\sqrt{-m})$. Auf diese Weise findet man sehr leicht, wenn
man aus den Formeln des vorigen Paragraphen $f(i)$, $f_1(\sqrt{-2})$,
$f(\sqrt{-3})$, $f(\sqrt{-7})$ entnimmt:
\weq{2}{f_1(\sqrt{-4})^8 = 8,}
\weq{3}{f_1(\sqrt{-16})^8 = 8\sqrt2(1+\sqrt2)^2,}
\weq{4}{f_1(\sqrt{-8})^8 = 8(1+\sqrt2),}
\weq{5}{f_1(\sqrt{-32})^8 = 32(1+\sqrt2)^2+8\sqrt2\sqrt{8(1+\sqrt2)^4+(1+\sqrt2)}.}
Setzen wir:
\[
f_1(\sqrt{-32})^8 = 8x,
\]
so ist $x$ Wurzel der biquadratischen Klassengleichung:
\weq{6}{(x^2-24x-2)^2-8(8x+1)^2=0,}
die durch Adjunktion von $\sqrt2$ in zwei quadratische Gleichungen
\[
x^2-8(1\pm\sqrt2)^2x-2(1\pm\sqrt2)=0
\]
zerfällt. Ferner finden wir so:
\weq{7}{f_1(\sqrt{-12})^4 = 2\sqrt[6]{2}(1+\sqrt3),}
\weq{8}{f_1(\sqrt{-28})^4 = 2\sqrt2(3+\sqrt7).}

Auf andere Fälle werden wir diese Methode später noch anwenden.

§ 130.

\gesperrt{Berechnung von Klasseninvarianten aus den Schlaeflischen Modulargleichungen.}

I. Wenn wir in einer der Gleichungen zwischen $u=f(\omega)$
und $v=f(n\omega)$ oder zwischen $u=f_1(\omega)$, $v=f_1(n\omega)$ (§ 73),
für $u$ einen der bekannten Werte von $f(\sqrt{-m})$ oder $f_1(\sqrt{-m})$
einsetzen, so erhalten wir eine Gleichung, welcher $f(\sqrt{-n^2m})$
oder $f_1(\sqrt{-n^2m})$ genügt.

Diese Gleichung enthält unter Umständen noch fremde Faktoren, die man
aufzusuchen und zu beseitigen hat.

Ist $n$ eine Primzahl, so erhalten wir nach § 123 vollständigen
Aufschluß über diese überflüssigen Faktoren. Geht $n$ in $m$ auf,
so ist in der betreffenden Gleichung ein zur Determinante $-m$
gehöriger Linearfaktor abzusondern, ist $-m$ quadratischer Nichtrest
von $n$, so sind überflüssige Faktoren überhaupt nicht vorhanden,
ist $-m$ quadratischer Rest von $n$, so ist ein zur Diskri-
%CONT>>
